\section{The determinant character on the prescribed higher theta level}
\label{sec:prescribed-theta-determinant}

Fix the original datum $D$, realization $R$, index $i$ and tuple $t$
of RTG1. Write $E/F$ for the host CM extension, $W=W_c(D,i)$,
$G=U(W)$, and $L= D.\mathrm{what}_2.\mathrm{LL}$.
Write $M=\operatorname{embLL}(L)\subset\mathbb C$ for the coefficient
field and $N=l_{2,t}.M$ for the integer level of RTG10.
We use the literal coefficient record with $c'.F=M$. The character
$c_t=R.\mathrm{xi}(t)$ is the adelic character tied by
$R.\mathrm{toLevCh.Is}(A.\mu_0)$ to the finite-order Hecke character
$R.\mathrm{xiH}(t)$. Choose its positive order bound $m_t$ from that
law. The original $\mathrm{XiVal}(R)$ says $c_t(x)\in M$ for every
idele $x$. In particular it makes an assertion about the values
of this particular character.

% BEGIN T73K1
\paragraph{T73K1. The ray character read on determinants.}
\label{claim:T73K1}
There is a homomorphism
$\xi_t:({\mathbb A}_{E,f})^\times\to\mathbb C^\times$
such that for $g\in G(\mathbb A_F)$ and every idele $x$ with
$x=\det(g)\bar x$,
\[
 \xi_t(\det g_f)=c_t(x)\in M,\qquad
 \xi_t(\det g_f)^{m_t}=1 .
 \tag{T73.1}
\]
It is trivial on determinants of rational unitary elements and on
determinants $u$ with $u-1\in N\widehat{\mathcal O}_E$.
One choice works for every adelic element and every root of its
determinant. This proves the higher-cover field $\mathrm{XiIs}$,
including the two triviality properties needed for descent.

\begin{proof}
The finite part of the Hecke character has exponent dividing $m_t$
by its order law. In the original idelic realization the $m_t$th power
is continuous, trivial on finite ideles and on diagonal $E^\times$.
By weak approximation $E^\times$ has dense image in the archimedean
factor, so this power is trivial everywhere. The archimedean factor
is a product of copies of $\mathbb C^\times$, hence is divisible,
and $c_t$ itself is trivial there. The base-triviality field of the
Hecke character now gives triviality on $\mathbb A_F^\times$.
Consequently $b/\bar b\mapsto c_t(b)$ on finite ideles is a
well-defined homomorphism: two lifts differ by a fixed idele, hence
by a finite idele of $F$. Its image has exponent dividing $m_t$.

Every norm-one finite idele admits such a lift. In a split factor
$(u,u^{-1})$ use $(u,1)$. In a quadratic field factor, Hilbert 90
gives a lift, explicitly $1+u$ if $u\ne-1$, or a nonzero trace-zero
element if $u=-1$. All nonsplit finite factors of this particular
host are unramified by $D.\mathrm{sp}.\mathrm{cm}.\mathrm{cmUnram}$.
There a uniformizer of $F_v$ is also a uniformizer of $E_w$, so the
lift may be multiplied by a fixed element to make it a unit. At
almost every split place $u$ is a unit. These choices define a finite
idele. In particular every unitary determinant lies in the domain
of the homomorphism just constructed. Extend it as an abstract
homomorphism to the full finite idele group using divisibility of
$\mathbb C^\times$: a maximal extension exists by unions of chains,
and on adjoining a new generator one chooses an $a$th root of the
already assigned value of its least positive power $a$ in the old
subgroup, or assigns $1$ if no such power exists. This gives the
precise domain required for $\xi_t$. Only its restriction to
determinants enters the assertions. A rational unitary determinant
admits a lift in $E^\times$ by the same Hilbert 90 calculation over
$E/F$, where the Hecke character is trivial. Formula (T73.1) follows
from base and archimedean triviality and $\mathrm{XiVal}$.

For the congruence assertion, $u\bar u=1$ and
$u-1\in N\widehat{\mathcal O}_E$ imply that both $u$ and $u^{-1}$
are integral. At every place over $N$, the level is odd, so
$b=(1+u)/2$ is a unit congruent to $1$ modulo $N$ and
$b=u\bar b$. Away from $N$ choose integral unit lifts as above.
The ramification law of $R$ places all ramified primes of its
character in $R.Q(t)$ or their conjugates. RTG10 gives $P\mid N$
for each such $P$, and conjugation fixes $N$. Hence away from $N$
the character kills units. At $P\mid N$ tameness kills the principal
units. The product of these local evaluations gives $c_t(b)=1$.
\end{proof}
% END T73K1

% BEGIN T73K2
\paragraph{T73K2. The component value is independent of every marking.}
\label{claim:T73K2}
Supply a higher model at the level of RTG10, its actual period family
$\mathcal F$, and a dictionary with the original $\mathrm{DictLat}$
for $\iota_W(T,i),x_0$, the $\mathrm{ModelW}$ identifications and the
original dictionary-domain condition. Write $j(h)=
\mathrm{ik}((\mathrm{dict}_2 h)_1)\in\mathrm{Fin}(\mathcal F.n)$.
For all witnesses $\mathrm{MarkBy}(h,\gamma)$, the value
$\xi_t(\det(\gamma h)_f)$ depends only on $j(h)$.

\begin{proof}
Let $g=\gamma h$, $g'=\gamma' h'$ be two marked representatives
with the same period index. The two lattice identifications transport
the identical alternating form of $\mathcal F.D_{j(h)}$ to
$\operatorname{Tr}_{E/\mathbb Q}(aH(v,w))$ on the two lattices,
with the same $a\in E^\times$ and the same Hermitian form $H$ by
$\mathrm{ModelW}$. Their composite extends uniquely to an
$E$-linear map $s$ between the rational spans, since both lattices are
full. For all lattice vectors $v,w$ and all $b\in\mathcal O_E$, insert
$bv$ in the trace equality. Nondegeneracy of the field trace pairing
and $a\ne0$ imply $H(sv,sw)=H(v,w)$, then the same holds throughout
the spans. Thus $s\in G(F)$.

The finite element $k=(sg)_f^{-1}g'_f$ preserves the original
completed lattice $\widehat\Lambda$, since the two lattices become
equal after $s$. The original $\mathrm{IotaOK}$ identifies the image
of $\iota_W(T,i)$ on $D_w.\mathrm{La}$ with the entire base lattice.
Completing this identification, both marks agree on
$\overline{\mathfrak b}_2\widehat\Lambda$ modulo
$N\widehat\Lambda$, so $k$ fixes that half of the marking.
Locally at a prime dividing $\mathfrak b_2$, the conjugate ideal is
a unit, whence $k\equiv 1$ modulo the local factor of $N$.
In an integral basis of the local lattice its determinant thus equals
$1$ modulo that factor. At primes dividing
$\overline{\mathfrak b}_2$ the conjugate prime lies in the first set.
Since $\det(k)\overline{\det(k)}=1$ the same determinant congruence
holds there as well. The two coprime half ideals multiply to $(N)$.
At all other primes lattice preservation makes the determinant a unit.
Therefore $\det(k)-1\in N\widehat{\mathcal O}_E$. By T73K1 both
$\det s$ and $\det k$ have character value one, which proves the
claim for arbitrary choices of both marks.
\end{proof}
% END T73K2

% BEGIN T73K3
\paragraph{T73K3. The algebraic character section over the original field.}
\label{claim:T73K3}
Keep the higher model and dictionary of T73K2. Give them their
original $\mathrm{IsUnf}$ and period-family laws, a higher
$\mathrm{DeltaDat}$ with $\mathrm{IsLev}$, and $\mathrm{CoverField}$
on the specified data. Then there exists a weight-zero, index-zero
higher section $f_t$ over $c_2.F=c'.F=M$ with
\[
 f_t^{m_t}=1,\qquad
 \mathrm{rdU}(U_2,f_t,\mathrm{dict}_2(h))
       =\xi_t(\det(\gamma h)_f)
 \quad\text{for every }\mathrm{MarkBy}(h,\gamma).
 \tag{T73.2}
\]
Substituting $f_t$ for the character entry of the given
$\mathrm{DeltaDat}$ preserves $\mathrm{IsLev}$. The scalar in
$\mathrm{CXiVal}$ can be $1$ simultaneously for every tuple.

\begin{proof}
For each period index $j$ occurring in a mark, put
$v_j=\xi_t(\det(\gamma h)_f)$ for any such mark. This is
well defined by T73K2. Put $v_j=1$ at all other indices. Every
$v_j$ is in $M$ and satisfies $v_j^{m_t}=1$ by T73K1.
The period index is locally constant on the complex points of
the higher space: the local period lifts keep the marked integral
lattice fixed, and the images for different marked classes are
disjoint by the equality-of-periods law. These are the disjoint
open marked-class images of HUF14 and their covering property.
The assertion also follows directly from the local lift and fiber
clauses of the given $\mathcal F$.

Apply $\mathrm{CoverField}$ to this particular level, higher model
and literal coefficient field. It supplies $\mathrm{FnAlgClosed}$.
For each irreducible component of the smooth finite-type scheme
$X_{o,2}\times_{K_A(l_2)}M$ this gives relative algebraic closedness
of $M$ in the component function field. The function field is a
regular extension, since $M$ has characteristic zero. Hence its
tensor product with $\mathbb C$ over the prescribed embedding is a
domain, and on affine opens the same holds by flatness. The components
are disjoint open and closed by smoothness. Each thus remains
irreducible after base change and has analytically connected complex
points (M.~Raynaud, SGA~1, expos\'e~XII, Proposition~2.4).
The integer $j$ is consequently constant there. For each $j$ let
$e_j$ be the indicator of the union of the corresponding components
over $M$. These are orthogonal algebraic idempotents with
$\sum_j e_j=1$.

The function $f_t=\sum_j v_j e_j$ gives a rule on every original
test ring by pulling it back via the underlying moduli point. Pullback
commutes with every test morphism, and triangular changes of frame and
a change of section generator leave this rule unchanged. These are the
two original weight-zero, index-zero section laws. On every test,
orthogonality gives $f_t^{m_t}=1$ as an equality of rules. At every
algebraically closed framed point the initial $\mathrm{IsLev}$ gives
a nonzero evaluation of its old section. The definition of
$\mathrm{evR}$ returns zero off the test locus, so each such point
belongs to the test locus; there evaluation is multiplicative and
takes $1$ to $1$. Therefore $\mathrm{evR}(f_t)^{m_t}=1$ there.
This proves the nonvanishing required by $\mathrm{IsLev}$ for the
replacement. Every other level and lifting field belongs to the same
given model. At an actual dictionary frame the dictionary-domain and
period-point laws identify its index with $j(h)$. Its reading is
$v_{j(h)}$, giving (T73.2) for every marking with the common scalar
$1$, of order one.
\end{proof}
% END T73K3

RTG12 can therefore use the actual determinant character and its
algebraic section once the level map, framed lifting, and marked higher
dictionary on the same $T,R$ are supplied. T73K1 is independent of
their construction. T73K2--T73K3 use the specified marked higher
geometric objects, the latter with $\mathrm{CoverField}$. The
coefficient comparison is the identity at $M$. Neither an enlargement
to split geometric components nor an extra cyclotomic coefficient
extension is part of this argument.
